\section{Introduction}\label{sec:intro}

Much company news is published before the stock market opens. Two questions follow. Can a model that reads that
news anticipate how the price will react? And once part of the reaction is visible, can it tell whether the
market moved too far or not far enough? Both are old questions in empirical finance. Prices react to firm news
within minutes \citep{patell1984,busse2002}, media tone predicts short-horizon returns \citep{tetlock2007}, and
moves that come with news tend to drift while moves without news tend to reverse \citep{chan2003}. Large
language models (LLMs) now read news at scale, and several studies find that their readings predict returns
\citep{lopezlira2026,chen2022}. We ask whether such predictability survives a test built to rule out the usual
ways a backtest misleads: look-ahead in the data, memorization in the model and selection over many variants.

We study \emph{AI decision models}. A decision model is a non-generative model that takes a \emph{state} (text
plus structured context) and a set of typed questions (a choice among options, a score on an ordered scale, or
a yes-or-no question) and returns probabilities over the answers instead of generating text
\citep{zhu2026}. Decision models are close relatives of LLM-based classifiers and of probabilistic
forecasting with language models \citep{kadavath2022,halawi2024,schoenegger2024}, and their vendors describe the
probabilities as calibrated. They suit finance because they return the quantities a forecaster needs, such as
a distribution over the size of the coming move or the probability that a piece of news is material, at a cost
of a few cents per thousand decisions, and because open-weights versions can be fine-tuned on free cloud GPUs.

Language models also bring a specific threat. A model trained on text published after the period on which it
is tested may know how events turned out, and even a model with an earlier training cutoff can recognize a
company from the content of its news \citep{glasserman2024,sarkar2024,lopezlira2025}. A backtest of a
language-model signal on historical news therefore mixes the model's ability to read news with its memory of
what followed. A second threat is older: when many signals, filters and horizons are tried, the best backtest
overstates what a strategy will earn \citep{harvey2016,bailey2014}. Pre-registration, fixing the analysis before
the data are seen, is the standard remedy in other sciences and is still rare in empirical finance
\citep{harvey2017,olken2015}.

We combine these elements in a pre-registered, look-ahead-free test. The sample is the point-in-time S\&P 500
from July 2024 to July 2026. For each company and trading session we build the news state as of 09:15~ET and
mask company names, tickers, people, dates and every mention of a price move. Before the open, each model gives
a \emph{pre-open reading}: how the stock will perform once the news is absorbed. At 10:00 and after the close
we show the model the reaction so far and ask for a \emph{post-reaction reading}: was the reaction too small or
too large for the news? The primary instrument, Laya-Fin, is an open-weights decision model that we fine-tuned
on 2024--2025 news and returns. We compare it with three commercial decision models chosen in a zero-shot
tournament, the untuned Laya model, an anchor model trained in the same way from an encoder released before
the sample, and standard baselines. Every definition, hypothesis and selection rule was written down before
any outcome was computed; selection used two validation quarters of 2025; the analysis plan was frozen with a
SHA-256 hash and a version-control tag before any 2026 model call; and the 2026 outcomes were unblinded once,
through a lock that logs every read. Every evaluated variant is logged as a trial, and every deviation from the
plan is logged with its date and reason. The study cost \$\res{spend-total} in API fees and \res{gpu-hours}
GPU-hours on a free cloud tier.

\needspace{8\baselineskip}
\paragraph{Key findings.}
\begin{enumerate}[label=(\arabic*), leftmargin=*, itemsep=3pt]
\item \textbf{Clean validation did not survive the blinded test.} The fine-tuned model's daily information
  coefficient (IC) with the close-to-close move was \res{lf-ic-vala} and \res{lf-ic-valb}
  ($t = \res{lf-ic-valb-t1}$) in the two 2025 validation quarters but \res{lf-ic-test}
  ($t = \res{lf-ic-test-t1}$) in the 2026 test; none of its five primary hypotheses is supported, and the
  commercial models' ICs roughly halved as well (d1-ultra from \res{d1-ic-valb} to \res{d1-ic-test}).
\item \textbf{Pre-open news is priced at the open.} The commercial decision models' daily IC with the overnight
  move from the previous close to the open is \res{d1-ic-on-test} to \res{drex-ic-on-test}
  ($t = \res{span-ic-on-test-t1}$ to \res{drex-ic-on-test-t1}), and their IC with the move after the open is
  about zero (\res{d1-ic-id-test} to \res{drex-ic-id-test}); no model predicts the move after the open.
\item \textbf{The models cannot tell over- from under-reaction.} In the pre-registered over/under regression
  the news coefficient is not positive (Laya-Fin $b = \res{lf-h2-b-d2}$, $p = \res{lf-h2-bp-d2}$) and the
  reversal coefficient has the predicted sign without significance ($c = \res{lf-h2-c-d2}$,
  $p = \res{lf-h2-cp-d2}$). Prices reach part of the model's news-implied value by the close (a slope of
  \res{lf-kappa-ct} for Laya-Fin and \res{d1-kappa-ct} for d1-ultra), then partly fade, and the gap's half-life
  is not reached within 20 trading days.
\item \textbf{Commercial decision models recognize ``anonymized'' firms.} Given masked news, they pick the
  company out of four same-sector candidates in \res{mem-min} to \res{mem-max} of cases, against a matched
  baseline of \res{mem-baseline}; the open-weights Laya model does so in \res{mem-laya}. This is a warning for
  backtests of language models on historical news.
\item \textbf{News explains little of a stock's own move out of sample.} Same-day news explains
  \res{lf-ladder-post-pct} of the stock-specific daily variance for Laya-Fin (supported; Benjamini--Hochberg
  $q = \res{lf-h5-q}$) and \res{d1-ladder-post-pct} for d1-ultra. Reading the realized move explains
  \res{d1-ladder-expl-pct}, but the moves the model calls explained do not persist. The models know which stocks
  will move a lot (an AUC of \res{lf-auc}), but no better than the stocks' own recent volatility (an
  out-of-sample $R^2$ gain of \res{lf-h4} over a realized-volatility benchmark).
\item \textbf{The commercial models are overconfident.} Laya-Fin is calibrated but compressed (expected
  calibration error \res{lf-ece}, log-loss skill \res{lf-llskill}), while the commercial models' stated
  probabilities overstate their hit rates (expected calibration error \res{drex-ece} to \res{span-ece}).
\item \textbf{No pre-registered strategy survives costs and selection.} No strategy clears its deflated Sharpe
  ratio (the largest probability is \res{dsr-max}, with \res{ntrials-strat} logged trials) or the 95th percentile
  of its random-signal placebo, and the convergence strategy is effectively a price-only gap fade
  (\res{sconv-shared} of positions shared, daily-return correlation \res{sconv-corr}).
\end{enumerate}

\paragraph{Contributions.} First, we give a protocol for testing language-model signals on financial news
without look-ahead: point-in-time states, masking checked by an independent scanner, a pre-registration frozen
by hash before any test-period call, a lock that releases test outcomes once, and ledgers of spend, trials and
test-data access, all in a replication package. Second, we show where the predictable part of pre-open news
goes: into the opening price. Third, we measure how prices converge to a model's news-implied value and read
the convergence path through an errors-in-variables lens that separates what its level and its shape
identify. Fourth, we add two measurement tools for language models in finance: a matched-distractor test of
whether a model recognizes masked companies, and a firm-level version of the explanation ladder of
\citet{kelly2026}, which measures how much of a stock's own daily move its news explains.

Section~\ref{sec:related} reviews related work and Section~\ref{sec:design} describes the data and the design.
Sections~\ref{sec:validation} to~\ref{sec:strategies} report the pre-registered results, nulls included.
Section~\ref{sec:exploratory} reports exploratory analyses, labelled post hoc, and Section~\ref{sec:limits} the
limitations and the pre-registered forward test.
